\textbf{UE :} S-INFO-015 Projet d'analyse et de conception logicielle\\
\textbf{Année académique :} 2026-2027\\
\textbf{Groupe :} G02 \textbf{Responsable de l'extension :} Vandooren
Romain (Matricule : 251074)

\subsubsection{Gestion des placements :}\label{gestion-des-placements}

\emph{Note :} Le terme « client investisseur » est défini comme un
client ayant un profil d'investisseur valide. Cette forme est utilisée
par simplicité d'écriture dans les User Stories ci-dessous.

\subparagraph{US1 - Demander l'ouverture d'un
compte-titre}\label{us1---demander-louverture-dun-compte-titre}

En tant que utilisateur connecté,\\
je veux ouvrir un compte titre\\
afin de réaliser des placements financiers.

\textbf{Critères d'acceptation :}

\begin{itemize}
\tightlist
\item
  Le compte-titre du client doit être lié à un de ses comptes courants
\item
  Ce compte-titre peut gérer plusieurs devises en allouant un
  compartiment pour chaque devise
\item
  Au sein d'un compartiment d'une , uniquement les placements utilisant
  cette devise y seront accessibles
\end{itemize}

\subparagraph{US2 - Remplir / modifier le formulaire de profil
d'investisseur}\label{us2---remplir-modifier-le-formulaire-de-profil-dinvestisseur}

En tant que client investisseur\\
je souhaite remplir un formulaire\\
afin de définir mon profil d'investisseur (conservateur, défensif,
neutre, dynamique, agressif)

\textbf{Critères d'acceptation :}

\begin{itemize}
\tightlist
\item
  A tout moment le formulaire peut être repassé par le client afin de
  modifier son profil
\item
  Le formulaire évalue la situation financière, le risque que le client
  est prêt à prendre et son horizon de placement (en années).
\item
  Le profil bloque tout nouvel achat de produit boursier dont le risque
  est plus élevé que celui configuré par le profil
\item
  Le formulaire doit être envoyé vers le serveur et être validé par
  l'institution financière
\item
  Le profil est attribué par l'institution financière
\end{itemize}

\subparagraph{US3 - Acheter des produits
financiers}\label{us3---acheter-des-produits-financiers}

En tant que client ayant un profil d'investisseur,\\
je veux acheter des produits financiers\\
afin d'enrichir mon compte-titre.

\textbf{Critères d'acceptation :}

\begin{itemize}
\tightlist
\item
  L'interface doit verrouiller l'achat et la visibilité des produits
  financiers dont le risque est strictement supérieur au profil du
  client
\item
  Les potentiels frais d'investissement et autres coûts liés sont inclus
  dans l'achat
\end{itemize}

\subparagraph{US4 - Consulter les placements réalisés du
compte-titre}\label{us4---consulter-les-placements-ruxe9alisuxe9s-du-compte-titre}

En tant que client investisseur,\\
je souhaite visualiser mes placements réalisés triés par type\\
afin de suivre l'évolution de mes investissements.

\textbf{Critères d'acceptation :}

\begin{itemize}
\tightlist
\item
  Pour chaque placement, l'interface affiche Nom, ISIN, Date achat, Date
  expiration (si existe), Quantité, Prix d'achat, Valeur actuelle,
  rendement, poids relatif dans le compte titre et dans le portefeuille
  global.
\end{itemize}

\subparagraph{US5 - Vendre des placements
financiers}\label{us5---vendre-des-placements-financiers}

En tant que client investisseur,\\
je souhaite vendre les placements de mon compte titre\\
afin d'obtenir de l'argent.

\textbf{Critères d'acceptation :}

\begin{itemize}
\tightlist
\item
  L'argent de la vente doit être placé sur un compte courant du client
\item
  Les taxes, frais de ventes et autres coûts liés doivent être inclus
  dans le processus de vente
\end{itemize}

\subparagraph{US6 - Comparer avec un portefeuille de
référence}\label{us6---comparer-avec-un-portefeuille-de-ruxe9fuxe9rence}

En tant que client investisseur\\
je souhaite voir un portefeuille de référence basé sur mon profil
d'investisseur\\
afin de comparer la performance et la diversification de mon
portefeuille actuel.

\textbf{Critères d'acceptation :}

\begin{itemize}
\tightlist
\item
  Un compte titre fictif doit être généré pour l'investisseur basé sur
  son profil de risque.
\item
  L'interface doit dédier une vue pour comparer le portefeuille réel et
  le portefeuille de référence.
\end{itemize}

\subparagraph{US7 - Consulter l'historique des prix et statistiques d'un
produit}\label{us7---consulter-lhistorique-des-prix-et-statistiques-dun-produit}

En tant que client investisseur,\\
je souhaite apercevoir l'historique des prix d'un produit à plusieurs
échelles de temps (D, W, M, 3M, 6M, YTD, 1Y, 5Y, MAX) sous forme d'un
graphique et de statistiques\\
afin de pouvoir réaliser un investissement éclairé.

\textbf{Critères d'acceptation :}

\begin{itemize}
\tightlist
\item
  Les données du produit boursier proviennent d'une API externe
\item
  Les données doivent être stocké dans le base de donnée
\item
  Les données doivent être mises à jour automatiquement régulièrement.
\end{itemize}

\subparagraph{US8 - Exporter les données du
compte-titre}\label{us8---exporter-les-donnuxe9es-du-compte-titre}

En tant que client investisseur,\\
je souhaite visualiser sous formes de graphique, exporter l'ensemble des
données ou une partie/période spécifique de mon portefeuille boursier\\
afin d'analyser mes performances en dehors de l'application.

\textbf{Critères d'acceptation :}

\begin{itemize}
\tightlist
\item
  L'interface présente des vues en tableaux de données et en graphiques.
\item
  Le client peut filtrer une période spécifique, une partie spécifique
  (par type de placement) ou regarder l'ensemble des données.
\item
  L'exportation peut se réaliser dans plusieurs formats différents (PDF,
  CSV)
\end{itemize}

\subparagraph{US9 - Percevoir les dividendes, intérêts et remboursements
à
terme}\label{us9---percevoir-les-dividendes-intuxe9ruxeats-et-remboursements-uxe0-terme}

En tant que client investisseur,\\
je souhaite percevoir automatiquement les dividendes de mes actions, les
intérêts de mes obligations et le remboursement de mes placements à
terme\\
afin de fructifier mon patrimoine financier.

\textbf{Critères d'acceptation :}

\begin{itemize}
\tightlist
\item
  Les dividendes et intérêts sont versés automatiquement sur le compte
  courant lié au compte titre à échéance ou à intervalle régulier selon
  le type de placement
\item
  Les prélèvement fiscaux liés aux revenus du capital sont
  automatiquement prélevés à la source si applicable
\item
  Lors de l'expiration d'un bon de caisse ou d'une obligation le montant
  investi est automatiquement reversé sur le compte courant lié au
  compte titre
\end{itemize}

\subparagraph{US10 - Valider et attribuer le profil
d'investisseur}\label{us10---valider-et-attribuer-le-profil-dinvestisseur}

En tant qu'employé d'une institution financière,\\
je souhaite analyser les formulaires soumis par mes clients\\
afin de leur attribuer un profil d'investisseur

\textbf{Critères d'acceptation :}

\begin{itemize}
\tightlist
\item
  L'employé a accès à chaque formulaire soumis par les clients de son
  institution financière
\item
  L'employé en fonction des réponses du formulaire soumis attribue un
  profil d'investisseur au client qui a soumis celui-ci
\item
  Une fois le formulaire validé, le formulaire soumis est archivé et le
  profil d'investisseur est appliqué au client
\end{itemize}

\subparagraph{US11 - Consulter les comptes-titres des clients de
l'institution}\label{us11---consulter-les-comptes-titres-des-clients-de-linstitution}

En tant qu'employé d'une institution financière,\\
je souhaite visualiser la liste des comptes-titres et leur contenu\\
afin de conseiller au mieux mes clients.

\textbf{Critères d'acceptation :}

\begin{itemize}
\tightlist
\item
  L'employé a uniquement accès aux comptes-titres des clients de
  l'institution financière où il travaille
\item
  L'interface de visualisation du compte-titre est identique à celle du
  client qui le détient
\item
  L'employé peut rechercher les comptes titres à l'aide de multiples
  critères (nom du titulaire, type de placement, compartiment de devise,
  etc\ldots)
\end{itemize}
